\documentclass[10pt,a4paper]{amsart}
\usepackage{amsmath,amssymb,amsthm,graphicx}
% Retained minimal preamble for the Stacks Project chapter flat.tex (More on Flatness).
% Verbatim excerpt of the chapter preamble: Xy-pic setup (native preamble lines 14-18) and the native
% theorem declarations with the mathematical macros (native preamble lines 154-217). Copied mechanically;
% nothing added, removed or redefined. The wrapper class and the amsmath/amssymb/amsthm/graphicx base
% packages are supplied by the assembler, so the class/externaldocument part of the native preamble is omitted.
% For commutative diagrams we use Xy-pic
\usepackage[all]{xy}

% We use 2cell for 2-commutative diagrams.
\xyoption{2cell}
% Theorem environments.
%
\theoremstyle{plain}
\newtheorem{theorem}[subsection]{Theorem}
\newtheorem{proposition}[subsection]{Proposition}
\newtheorem{lemma}[subsection]{Lemma}

\theoremstyle{definition}
\newtheorem{definition}[subsection]{Definition}
\newtheorem{example}[subsection]{Example}
\newtheorem{exercise}[subsection]{Exercise}
\newtheorem{situation}[subsection]{Situation}

\theoremstyle{remark}
\newtheorem{remark}[subsection]{Remark}
\newtheorem{remarks}[subsection]{Remarks}

\numberwithin{equation}{subsection}

% Macros
%
\def\lim{\mathop{\mathrm{lim}}\nolimits}
\def\colim{\mathop{\mathrm{colim}}\nolimits}
\def\Spec{\mathop{\mathrm{Spec}}}
\def\Hom{\mathop{\mathrm{Hom}}\nolimits}
\def\Ext{\mathop{\mathrm{Ext}}\nolimits}
\def\SheafHom{\mathop{\mathcal{H}\!\mathit{om}}\nolimits}
\def\SheafExt{\mathop{\mathcal{E}\!\mathit{xt}}\nolimits}
\def\Sch{\mathit{Sch}}
\def\Mor{\mathop{\mathrm{Mor}}\nolimits}
\def\Ob{\mathop{\mathrm{Ob}}\nolimits}
\def\Sh{\mathop{\mathit{Sh}}\nolimits}
\def\NL{\mathop{N\!L}\nolimits}
\def\CH{\mathop{\mathrm{CH}}\nolimits}
\def\proetale{{pro\text{-}\acute{e}tale}}
\def\etale{{\acute{e}tale}}
\def\QCoh{\mathit{QCoh}}
\def\Ker{\mathop{\mathrm{Ker}}}
\def\Im{\mathop{\mathrm{Im}}}
\def\Coker{\mathop{\mathrm{Coker}}}
\def\Coim{\mathop{\mathrm{Coim}}}

% Boxtimes
%
\DeclareMathSymbol{\boxtimes}{\mathbin}{AMSa}{"02}

%
% Macros for moduli stacks/spaces
%
\def\QCohstack{\mathcal{QC}\!\mathit{oh}}
\def\Cohstack{\mathcal{C}\!\mathit{oh}}
\def\Spacesstack{\mathcal{S}\!\mathit{paces}}
\def\Quotfunctor{\mathrm{Quot}}
\def\Hilbfunctor{\mathrm{Hilb}}
\def\Curvesstack{\mathcal{C}\!\mathit{urves}}
\def\Polarizedstack{\mathcal{P}\!\mathit{olarized}}
\def\Complexesstack{\mathcal{C}\!\mathit{omplexes}}
% \Pic is the operator that assigns to X its picard group, usage \Pic(X)
% \Picardstack_{X/B} denotes the Picard stack of X over B
% \Picardfunctor_{X/B} denotes the Picard functor of X over B
\def\Pic{\mathop{\mathrm{Pic}}\nolimits}
\def\Picardstack{\mathcal{P}\!\mathit{ic}}
\def\Picardfunctor{\mathrm{Pic}}
\def\Deformationcategory{\mathcal{D}\!\mathit{ef}}

% --- appended 2026-10-12 (Chief): the project own definitions of the reference/slogan/history
% display environments, copied mechanically from the shared Stacks preamble lines 8-12
% (\usepackage{verbatim} plus the three \newenvironment lines, which the chapter preamble keeps
% verbatim from the project preamble). Nothing added, removed or redefined: these environments
% carry no mathematical content and the project itself typesets them as comments. Without these
% three lines a counted statement whose own text uses them fails with Environment slogan undefined.
\usepackage{verbatim}
\newenvironment{reference}{\comment}{\endcomment}
\newenvironment{slogan}{\comment}{\endcomment}
\newenvironment{history}{\comment}{\endcomment}
\usepackage{cleveref}
\crefname{originallemma}{original lemma}{original lemmas}
\Crefname{originallemma}{Original lemma}{Original lemmas}
\makeatletter
\newlabel{lemma-flat-finite-type-finitely-presented-over-dense-open-X}{{081P}{source}{}{sourcepointer-1}{}}
\newlabel{lemma-flat-finite-type-finitely-presented-over-dense-open-X@cref}{{[originallemma][0][]081P}{[1][0][]source}{}{}{}}
\makeatother
\crefname{originallemma}{original lemma}{original lemmas}
\Crefname{originallemma}{Original lemma}{Original lemmas}
\makeatletter
\newlabel{more-morphisms-lemma-flat-finite-presentation-CM-open}{{045U}{source}{}{sourcepointer-2}{}}
\newlabel{more-morphisms-lemma-flat-finite-presentation-CM-open@cref}{{[originallemma][0][]045U}{[1][0][]source}{}{}{}}
\makeatother
\crefname{originallemma}{original lemma}{original lemmas}
\Crefname{originallemma}{Original lemma}{Original lemmas}
\makeatletter
\newlabel{more-morphisms-lemma-flat-finite-presentation-CM-pieces}{{054T}{source}{}{sourcepointer-3}{}}
\newlabel{more-morphisms-lemma-flat-finite-presentation-CM-pieces@cref}{{[originallemma][0][]054T}{[1][0][]source}{}{}{}}
\makeatother
\crefname{originaldefinition}{original definition}{original definitions}
\Crefname{originaldefinition}{Original definition}{Original definitions}
\makeatletter
\newlabel{morphisms-definition-relative-dimension-d}{{02NJ}{source}{}{sourcepointer-4}{}}
\newlabel{morphisms-definition-relative-dimension-d@cref}{{[originaldefinition][0][]02NJ}{[1][0][]source}{}{}{}}
\makeatother
\crefname{originallemma}{original lemma}{original lemmas}
\Crefname{originallemma}{Original lemma}{Original lemmas}
\makeatletter
\newlabel{morphisms-lemma-fppf-open}{{01UA}{source}{}{sourcepointer-5}{}}
\newlabel{morphisms-lemma-fppf-open@cref}{{[originallemma][0][]01UA}{[1][0][]source}{}{}{}}
\makeatother
\title{Bounding the relative dimension from a dense open, for a flat morphism}
\author{The Stacks Project Authors}
\date{stacks-project:flat.tex:local-snapshot-20261011; GNU FDL 1.2 or later}
\begin{document}
\maketitle
\noindent\textit{Editorial note.} All delivered mathematical text is the author's own, copied from the pinned Stacks Project chapter file flat.tex; no statement or proof step was rewritten, reordered or removed.\par\medskip
\noindent\textit{Editorial note.} Preamble provenance: the chapter's own retained minimal preamble (sources/stacks-flat/retained-preamble.tex, lines 1-85) is a verbatim excerpt of the native Xy-pic setup and of the native theorem declarations and mathematical macros, followed by the project's own reference/slogan/history display-environment definitions copied mechanically from the shared preamble; the wrapper class and the amsmath/amssymb/amsthm/graphicx base packages are supplied by the assembler, and extra\_packages is left empty because everything needed is already in the retained preamble. No macro was invented or redefined.\par\medskip
\noindent\textit{Editorial note.} Same-chapter pointers are located at their \textbackslash{}label lines inside the pinned chapter file; cross-chapter pointers are located in the project's own tag list staged in the snapshot (sources/stacks-chow/tags.txt), at the verbatim lines carrying their canonical tags.\par\medskip
\noindent\textit{Allowed standard prerequisites (editorial).}\begin{itemize}
\item The open subscheme of the source on which the morphism has bounded fibre dimension, and its decomposition into pieces. More on Morphisms, labels more-morphisms-lemma-flat-finite-presentation-CM-open and more-morphisms-lemma-flat-finite-presentation-CM-pieces, invoked at proof lines 3035-3036. Provides the open subscheme containing every generic point of every fibre and its decomposition by the fibre dimension, which reduces the statement to the vanishing of the pieces of dimension larger than the bound.
\item Openness of a flat locally finitely presented morphism, and the dense open case with retrocompactness. Morphisms, label morphisms-lemma-fppf-open (proof line 3042); Same chapter flat.tex, label lemma-flat-finite-type-finitely-presented-over-dense-open-X (proof line 3049). The first makes the image of each piece open so that it must meet the given dense open, and the second turns the retrocompact dense open case into the locally finitely presented case.
\end{itemize}\medskip
\ifdefined\hypertarget\hypertarget{sourcepointer-1}{}\fi
\noindent\textit{Source reference.} Original lemma 081P: sources/stacks-flat/pinned-original-source.tex, source lines 3000--3000. Navigation pointer for \textbackslash{}ref\{lemma-flat-finite-type-finitely-presented-over-dense-open-X\} used inside the counted span (line 3049); the argument that consumes it is retained in the span.\par\medskip
\ifdefined\hypertarget\hypertarget{sourcepointer-2}{}\fi
\noindent\textit{Source reference.} Original lemma 045U: sources/stacks-chow/tags.txt, source lines 5449--5449. Navigation pointer for \textbackslash{}ref\{more-morphisms-lemma-flat-finite-presentation-CM-open\} used inside the counted span (line 3036); the argument that consumes it is retained in the span.\par\medskip
\ifdefined\hypertarget\hypertarget{sourcepointer-3}{}\fi
\noindent\textit{Source reference.} Original lemma 054T: sources/stacks-chow/tags.txt, source lines 6651--6651. Navigation pointer for \textbackslash{}ref\{more-morphisms-lemma-flat-finite-presentation-CM-pieces\} used inside the counted span (line 3037); the argument that consumes it is retained in the span.\par\medskip
\ifdefined\hypertarget\hypertarget{sourcepointer-4}{}\fi
\noindent\textit{Source reference.} Original definition 02NJ: sources/stacks-chow/tags.txt, source lines 3482--3482. Navigation pointer for \textbackslash{}ref\{morphisms-definition-relative-dimension-d\} used inside the counted span (line 3024); the argument that consumes it is retained in the span.\par\medskip
\ifdefined\hypertarget\hypertarget{sourcepointer-5}{}\fi
\noindent\textit{Source reference.} Original lemma 01UA: sources/stacks-chow/tags.txt, source lines 2403--2403. Navigation pointer for \textbackslash{}ref\{morphisms-lemma-fppf-open\} used inside the counted span (line 3042); the argument that consumes it is retained in the span.\par\medskip
\begin{lemma}
\label{lemma-flat-finite-presentation-dimension-over-dense-open}
Let $f : X \to S$ be a morphism of schemes which is flat and locally
of finite type. Let $U \subset S$ be a dense open such that
$X_U \to U$ has relative dimension $\leq e$, see
Morphisms, Definition \ref{morphisms-definition-relative-dimension-d}.
If also either
\begin{enumerate}
\item $f$ is locally of finite presentation, or
\item $U \subset S$ is retrocompact,
\end{enumerate}
then $f$ has relative dimension $\leq e$.
\end{lemma}

\begin{proof}
Proof in case (1). Let $W \subset X$ be the open subscheme constructed
and studied in More on Morphisms, Lemmas
\ref{more-morphisms-lemma-flat-finite-presentation-CM-open} and
\ref{more-morphisms-lemma-flat-finite-presentation-CM-pieces}.
Note that every generic point of every fibre is contained in $W$,
hence it suffices to prove the result for $W$. Since
$W = \bigcup_{d \geq 0} U_d$, it suffices to prove that $U_d = \emptyset$
for $d > e$. Since $f$ is flat and locally of finite presentation it is open
hence $f(U_d)$ is open (Morphisms, Lemma \ref{morphisms-lemma-fppf-open}).
Thus if $U_d$ is not empty, then $f(U_d) \cap U \not = \emptyset$ as
desired.

\medskip\noindent
Proof in case (2). We may replace $S$ by its reduction. Then $U$ is
scheme theoretically dense. Hence $f$ is locally of finite presentation
by Lemma \ref{lemma-flat-finite-type-finitely-presented-over-dense-open-X}.
In this way we reduce to case (1).
\end{proof}

\end{document}
